% !TEX root = lectures.tex
\section{Lecture 2}
To extend our results from before
to the case where there are more than two possible predictions,
we can think of ``following'' an expert instead of
picking a specific option.
We can model $\ell_t(i)$ as the ``loss'' (error measure) of expert $i$
at time $t$. We can generalize the loss to model fractional errors, allowing $\ell_t(i) \in [0, \infty)$.

\begin{algothm}[Hedge Algorithm]
    We again have weights $\{W_t(i)\}$ for the experts.
    We pick advice $i$ time $t$ with probability $\frac{W_t(i)}{\sum_j W_t(j)} =: p_t(i)$.
    We suffer (expected) loss $\E{L} = \sum_i p_t(i) \ell_t(i)$. Then,
    we update weights $W_{t + 1}(i) = W_t(i) e^{-\epsilon \ell_t(i)}$.
\end{algothm}

Our yardstick is pretty similar.

\begin{theorem}
    For all $i$,
    \[ \sum_{t = 1}^T p_t^{\top} \ell_t \le \sum_{t=1}^T \ell_t(i) + \epsilon \sum_{t=1}^{T} p_t^{\top} \ell_t^2 + \frac{\log N}{\epsilon}  \] 
    \begin{proof*}
        We again define $\phi_t = \sum_i W_{t}(i)$.
        \begin{align*}
            \phi_{t + 1} &= \sum_{i = 1}^{N} W_t(i) e^{-\epsilon \ell_t(i)} \\
            &= \phi_t \sum_{i = 1}^N p_t(i) e^{-\epsilon \ell_t(i)} \\
            &\le \phi_t \sum_{i = 1}^N p_t(i) (1 - \epsilon \ell_t(i) + \epsilon^2 \ell^2_t(i)) \\
            &\le \phi_t e^{-\epsilon p_t^{\top} \ell_t + \epsilon^2 p_t^{\top} \ell_t^2 }
        \end{align*}
        Furthermore, the potential function is at least $\Phi_T \ge e^{-\epsilon \sum_t \ell_t(i)}$.
        Now $\Phi_T \le N e^{-\epsilon p_t^{\top} \ell_t + \epsilon^2 p_t^{\top} \ell_t^2 }$.
        Taking a log and dividing by $\epsilon$ gives the result.
    \end{proof*}
\end{theorem}

Observe that if $\ell_t(i) \in [0, 1]$, then the expected loss $\E{L} \le \sum_t \ell_t(i) + \epsilon T + \frac{\log N}{\epsilon}$.
If we pick $\epsilon = \sqrt{\frac{\log N}{T}}$, then
\[ \frac{1}{T} \sum_t \qty(p_t^{\top} \ell_t - \ell_t(i)) \le 2 \sqrt{\frac{\log N}{T}} \to 0\]
We can define regret as $R = \sum_t p_t^{\top} \ell_t - \ell_t(i)$.

\subsection{Relationship to Other ML Problems}
Consider the \textbf{online routing} problem. There is some graph $G = (V, E)$
with a source $s$ and destination $t$. We wish to find the fastest path from $s$ to $t$.
The time taken to traverse each path constantly changes due to network traffic.
We can reduce this back to the experts problems, where each path is an expert.
There can be as many as $N!$ paths, so the regret will be $R \sim \sqrt{|E| \frac{\log N!}{T}} \sim \sqrt{\frac{N^3 \log N}{T}}$.

Consider the \textbf{spam classification} problem.
Represent an email with a bag-of-words representation $v \in \{0, 1\}^{d}$.
where $d = |\text{dictionary}|$. Then, a classifier
is a vector $x \in \R^d$ with $|x| \le 1$ and uses $\text{sign}(x^{\top} v)$
to report spam or not spam. We can take ``all possible''
classifiers up to resolution $\epsilon$ where
there are $\qty(\frac{1}{\epsilon})^{d}$ such classifiers. These
will form our expert set, and then we can just run hedge.

However, both of these approaches miss out on the structure of the underlying problems.
For instance, there are only $|E|$ numbers that determine the shortest path in the graph,
but we seem to elide the fact that the paths largely share edges with each other.

\subsection{Exploiting Structure}
Structured decisions lead us to the world of mathematical optimization.
\begin{enumerate}
\item The set of all norm at most $1$ linear classifiers is
just the unit ball.
\item The set of all $s$-$t$ paths in $G = (V, E)$ form a simplex.
In particular, for some path $P_{st}$, define $x_e = 1$ if $e \in P_{st}$ 
and $0$ otherwise. Then define set $K \subseteq R^{|E|}$,
\[ K = \left\{ x\in \R^{|E|}, \sum_{i \in V} x_{si} = 1, \sum_{i \in V} x_{it} = 1, \forall j \sum_i x_{ij} = \sum_{k} x_{jk}, x_{ij} \in [0, 1]\right\} \]
One may recognize this as the flow polytope. The corners of this polytope form paths from $s$ to $t$.
\end{enumerate}